% icstdoc.tex V1.12 6 May 2011

\documentclass[fonts]{icst}

\usepackage{moreverb}

\usepackage[breaklinks,colorlinks,bookmarksopen,bookmarksnumbered,linkcolor=ICSTblue,citecolor=blue,urlcolor=ICSTblue]{hyperref}
\usepackage{breakurl}
\usepackage{doi}
\usepackage{subfiles}
\usepackage{graphicx}
\usepackage{amsmath}
\usepackage{amssymb}
\usepackage{booktabs}
\usepackage{hyperref}
\usepackage{url}
\usepackage{svg}
\usepackage{algorithm}
\usepackage{algorithmic}
\usepackage{multirow}
\usepackage{xcolor}
\usepackage{tikz}
\usetikzlibrary{arrows.meta, positioning, fit, backgrounds, calc}

\newcommand\BibTeX{{\rmfamily B\kern-.05em \textsc{i\kern-.025em b}\kern-.08em
T\kern-.1667em\lower.7ex\hbox{E}\kern-.125emX}}

\def\volumeyear{2026}

\journalname{EAI Endorsed Transactions on AI and Robotics}
\articletype{Research Article}
\setcounter{page}{01}

  \def\copyrightholder{Q.-V. Dang\textit{ et al.}, licensed to EAI}
  \copyrightnote{This is an open access article distributed under the terms of the \href{https://creativecommons.org/licenses/by-nc-sa/4.0/}{CC BY-NC-SA 4.0}, which permits copying, redistributing, remixing, transformation, and building upon the material in any medium so long as the original work is properly cited.}
  \def\articledoi{10.4108/XX.X.X.XX}
\received{XXXX}
  \accepted{XXXX}
  \published{XXXX}

\begin{document}

\runningheads{Q.-V. Dang \textit{et al.}}{XHBot: eXplainable Heterophily-aware Graph Neural Networks for Social Bot Detection}

\title{XHBot: eXplainable Heterophily-aware Graph Neural Networks for Social Bot Detection}

\author{Quang-Vinh Dang\affil{1}, Phuong-Lan Nguyen\affil{1}, Dat Le\affil{2}\fnoteref{1}, Minh Ngoc Dinh\affil{3}}

\address{\affilnum{1}British University Vietnam, Hung Yen, Vietnam\\
\affilnum{2}University of Economics and Finance, Ho Chi Minh City, Vietnam\\
\affilnum{3}Millennia Education, Ho Chi Minh City, Vietnam}

\abstract{
Social bots threaten the integrity of online ecosystems by engaging in coordinated opinion manipulation. While Graph Neural Networks (GNNs) have become a dominant paradigm for bot detection, modern camouflaged bots strategically follow benign users to evade detection, creating structural heterophily that degrades the performance of standard homophilic GNN aggregators; moreover, many existing detectors offer limited forensic explainability. To address these challenges jointly, we propose XHBot (eXplainable Heterophily-aware Bot detector), a framework that is robust to heterophilic relation camouflage while providing transparent, multi-level forensic evidence for platform moderation. XHBot couples three components: Spectral-Guided Topology Refinement (SGTR), which down-weights camouflage edges by their contribution to the graph's high-frequency (Dirichlet) energy before aggregation; Tri-Channel Heterophily-Aware Aggregation (THCA), which separates homophilic, heterophilic, and self-identity signals; and Contrastive Prototype Disentanglement (CPD), which decouples behavioural signatures from social positioning. Evaluated on TwiBot-20, TwiBot-22, and Cresci-2017 under a unified protocol, XHBot reaches an F1 score of 0.9474 on TwiBot-20, improving over a competitive suite of recent baselines (including RGT, NeighborSense, and HW-GNN) by 9.64\%. Its Hierarchical Forensic Explanation (HFE) module extracts both instance-level subgraphs and community-level diagnostic motifs, which we assess quantitatively (Fidelity, Sparsity) and through qualitative case studies. These results indicate that decoupling behavioural signatures from adversarial social positioning is valuable for modern bot detection, and that combining accuracy with interpretable evidence supports deployment in real-world moderation settings.
}

\keywords{Social Bot Detection, Graph Neural Networks, Heterophily, Explainable AI, \emph{XHBot}}

\fnotetext[1]{Corresponding author.  Email: \email{datla@uef.edu.vn}}

\maketitle

% \bibliography{ref.bib}
\subfile{1_intro}
\subfile{2_related_work}
\subfile{3_method}
\subfile{4_experiment}
\subfile{5_discussion}
\subfile{6_conclusion}

\bibliographystyle{icstnum}
\bibliography{ref}
\end{document}
